\begin{tikzpicture}[font=\small,line cap=round,line join=round]\begin{scope}[x=0.9cm,y=0.19047619047619047cm]\draw[otRule!65,line width=.2pt](-5,-1)--(-5,20);\node[below,font=\tiny,fill=white,inner sep=.5pt]at(-5,0){$-5$};\draw[otRule!65,line width=.2pt](-4,-1)--(-4,20);\node[below,font=\tiny,fill=white,inner sep=.5pt]at(-4,0){$-4$};\draw[otRule!65,line width=.2pt](-3,-1)--(-3,20);\node[below,font=\tiny,fill=white,inner sep=.5pt]at(-3,0){$-3$};\draw[otRule!65,line width=.2pt](-2,-1)--(-2,20);\node[below,font=\tiny,fill=white,inner sep=.5pt]at(-2,0){$-2$};\draw[otRule!65,line width=.2pt](-1,-1)--(-1,20);\node[below,font=\tiny,fill=white,inner sep=.5pt]at(-1,0){$-1$};\draw[otRule!65,line width=.2pt](0,-1)--(0,20);\draw[otRule!65,line width=.2pt](1,-1)--(1,20);\node[below,font=\tiny,fill=white,inner sep=.5pt]at(1,0){$1$};\draw[otRule!65,line width=.2pt](2,-1)--(2,20);\node[below,font=\tiny,fill=white,inner sep=.5pt]at(2,0){$2$};\draw[otRule!65,line width=.2pt](3,-1)--(3,20);\node[below,font=\tiny,fill=white,inner sep=.5pt]at(3,0){$3$};\draw[otRule!65,line width=.2pt](4,-1)--(4,20);\node[below,font=\tiny,fill=white,inner sep=.5pt]at(4,0){$4$};\draw[otRule!65,line width=.2pt](5,-1)--(5,20);\node[below,font=\tiny,fill=white,inner sep=.5pt]at(5,0){$5$};\draw[otRule!65,line width=.2pt](-5,0)--(5,0);\draw[otRule!65,line width=.2pt](-5,5)--(5,5);\node[left,font=\tiny,fill=white,inner sep=.5pt]at(0,5){$5$};\draw[otRule!65,line width=.2pt](-5,10)--(5,10);\node[left,font=\tiny,fill=white,inner sep=.5pt]at(0,10){$10$};\draw[otRule!65,line width=.2pt](-5,15)--(5,15);\node[left,font=\tiny,fill=white,inner sep=.5pt]at(0,15){$15$};\draw[otRule!65,line width=.2pt](-5,20)--(5,20);\node[left,font=\tiny,fill=white,inner sep=.5pt]at(0,20){$20$};\draw[otMuted,->](-5,0)--(5.166666666666667,0)node[right]{$x$};\draw[otMuted,->](0,-1)--(0,21.05)node[above]{$y$};\node[below left,font=\tiny]at(0,0){0};\begin{scope}\clip(-5,-1)rectangle(5,20);\draw[otNavy,line width=1pt]plot[domain=-4.5:4.5,samples=90](\x,{\x*\x});\geoSolution{\draw[geoAide](-5,16)--(5,16);}\end{scope}\end{scope}\end{tikzpicture}
{\par\smallskip\fontsize{8}{10}\selectfont\color{otMuted}$y=x^2$. Les axes sont gradués ; leurs échelles peuvent différer.\par}
